\section{Retrieval + Reasoning: primero encontrar la estructura relevante y después desarrollar el cálculo}

Algunas fallas no se deben a que el modelo no sepa razonar, sino a que a la entrada le faltan hechos, fórmulas, precedentes o abstracciones de nivel superior. Retrieval y reasoning resuelven etapas distintas: el primero cambia la información que el modelo puede ver; el segundo ejecuta la descomposición y el cálculo sobre la información dada. En clase se rechaza la disputa de «elegir uno de los dos» y se propone directamente retrieval + reasoning, porque el objetivo del producto es la exactitud global y no defender un bando metodológico.

\subsection{Por qué la recuperación y el razonamiento no se sustituyen entre sí}

La pregunta de la slide 43 ocupa una sola línea, pero resume la división del sistema en capas: el recuperador decide qué evidencia entra en el contexto y el razonador decide cómo combinar esa evidencia para llegar a una conclusión. Si la recuperación acierta pero el razonamiento falla, aparecen citas correctas con conclusiones erróneas; si el razonamiento es correcto pero falta la recuperación, se calcula con rigor sobre premisas equivocadas. Ambos deben evaluarse por separado en recall, provenance, entailment y éxito final de la tarea.

\lecturefigure{slide-43.jpg}{¿Retrieval or reasoning? La respuesta de la clase es Retrieval + Reasoning}{43}

Supongamos que el recuperador devuelve, a partir del corpus $\mathcal C$, $z\sim q_\phi(z\mid x)$, y que el generador produce después $y\sim p_\theta(y\mid x,z)$. La probabilidad ideal de la respuesta es
\[
p(y\mid x)=\sum_{z\in\mathcal C}q_\phi(z\mid x)p_\theta(y\mid x,z).
\]
Aquí $\phi$ son los parámetros de recuperación, $\theta$ los parámetros de generación y $z$ la evidencia. Los sistemas reales solo toman los documentos top-$k$ para aproximar la suma; por eso $k$, el reordenador, la compresión del contexto y la conservación de las citas modifican el resultado. La recuperación no termina con meter el texto de una página web en el prompt: es un componente probabilístico con tasa de recuperación, ruido y vigencia temporal.

\begin{knowledgebox}{Primer uso: Retrieval-Augmented Reasoning}
Consiste en obtener primero información relevante de documentos externos, memoria, herramientas o bancos de casos, y después hacer que el modelo genere los pasos intermedios y la respuesta con base en esa evidencia. Es indispensable conservar la fuente y la marca de tiempo, y distinguir entre «el modelo lo recuerda», «se recuperó» y «se dedujo de la evidencia»; de lo contrario, no es posible localizar los errores.
\end{knowledgebox}

\subsection{Analogical Reasoning: recuperar problemas relacionados, no la misma respuesta}

La slide 44 presenta un cuadrado rotado cuyos cuatro vértices son $(-2,2)$, $(2,-2)$, $(-2,-6)$ y $(-6,-2)$. Si se pregunta directamente por el área, el modelo puede fallar; tras la indicación «recall a related problem», primero evoca la fórmula de la distancia entre dos puntos. La distancia entre vértices adyacentes es
\[
s=\sqrt{(2-(-2))^2+(-2-2)^2}=\sqrt{32},
\qquad A=s^2=32.
\]
Aquí $s$ es la longitud del lado del cuadrado y $A$ su área. Lo que se recupera no es el mismo problema, sino una subestructura transferible, «la distancia entre dos puntos»; esto muestra que retrieval puede devolver primitivas algorítmicas, plantillas de demostración o casos de falla similares, y no solo párrafos con hechos.

\lecturefigure{slide-44.jpg}{Analogical reasoning: primero se evoca un problema de distancia relacionado y después se calcula el área del cuadrado rotado}{44}

\teachervoice{00:51:29--00:55:00, el docente dice que en la industria le importa más el desempeño y que no quiere discutir si es retrieval o reasoning. Nota de clase: un related problem no equivale a un duplicate problem; una recuperación útil suele aportar principios transferibles, no filtrar directamente la respuesta.}

\begin{warningbox}{Transferencia negativa en la recuperación por analogía}
Los casos parecidos en la superficie pueden tener restricciones, unidades o mecanismos causales distintos. El sistema debe hacer que el modelo escriba explícitamente «qué estructuras son iguales y qué condiciones difieren», y verificar las premisas antes de aplicar la analogía; de lo contrario, la recuperación por similitud introducirá en el razonamiento la plantilla equivocada que más se parece.
\end{warningbox}

\subsection{Step-Back: primero abstraer, luego volver a los detalles}

El step-back prompting de la slide 45 no consiste simplemente en «pensar más tiempo», sino en elevar el problema original a un concepto o principio más general: en un problema de física se busca primero la ley de conservación; en preguntas con restricciones temporales se ordena primero la trayectoria completa y luego se vuelve al intervalo concreto. Esta operación reescribe la query de recuperación, de la redacción superficial a una pregunta estructural, y con frecuencia permite encontrar conocimiento más estable.

\lecturefigure{slide-45.jpg}{Step-Back Prompting: evocar principios mediante la abstracción y después resolver el problema original}{45}

El flujo puede escribirse en tres pasos:
\begin{enumerate}
  \item Generar la pregunta abstracta $x' = g(x)$, por ejemplo «¿qué relación de conservación sigue este tipo de sistema?»;
  \item Recuperar o generar el principio $z=\operatorname{retrieve}(x')$;
  \item Resolver bajo las restricciones originales $y=\operatorname{reason}(x,z)$.
\end{enumerate}
El criterio de aceptación decisivo es si $z$ realmente se aplica a $x$. Si la abstracción es demasiado amplia, la recuperación devuelve obviedades; si es demasiado estrecha, degenera en una reformulación del problema original. Conviene registrar la query abstracta, los aciertos de evidencia y las citas finales, para que el análisis de fallas sepa si el problema estuvo en el query rewrite, en retrieval o en reasoning.

\subsection{Deep Research es un ciclo de varias rondas de recuperación y razonamiento}

La slide 46 usa el acceso al producto Deep Research para mostrar cómo la misma idea se extiende a tareas largas: el sistema no recupera una sola vez, sino que formula subpreguntas, busca, lee, actualiza hipótesis, sigue recuperando y, al final, organiza un informe con fuentes. En la pantalla de la clase aparecen a la vez el título Gemini Deep Research y los accesos de producto de búsqueda/Deep research; el énfasis didáctico está en la categoría de producto, y la UI no debe tomarse como garantía de funciones en 2026.

\lecturefigure{slide-46.jpg}{Deep Research: organizar la recuperación, la descomposición y la síntesis como una tarea de varias rondas}{46}

\begin{lstlisting}[language=Python,caption={Ciclo de retrieval-reasoning con provenance y condición de paro}]
def deep_research(question, search, reasoner, budget):
    state = {"open": [question], "evidence": [], "claims": []}
    while state["open"] and budget.remaining():
        subquestion = state["open"].pop(0)
        documents = search(subquestion, top_k=8)
        update = reasoner.synthesize(subquestion, documents)
        state["evidence"].extend(update.citations)
        state["claims"].extend(update.claims)
        state["open"].extend(update.followups)
        budget.charge(documents, update)
    return reasoner.final_report(state, require_citations=True)
\end{lstlisting}

Este ciclo necesita una condición de paro y la eliminación de evidencia duplicada. Sin un límite de presupuesto, el sistema puede seguir planteando preguntas nuevas indefinidamente; sin provenance, el texto final no puede verificarse; sin alineación claim--evidence, las citas son solo decorativas. Para los hechos que cambian con el tiempo, también hay que guardar la fecha de obtención, porque el contenido de una misma URL puede cambiar.

\teachervoice{00:55:00--00:56:31, el docente vincula analogical reasoning, step-back y deep research: primero encontrar el conocimiento relevante o una estructura de nivel superior, y después razonar. Las notas advierten que un deep research real también debe manejar la calidad de las fuentes, los conflictos, la vigencia y la terminación; no se vuelve confiable automáticamente por buscar unas cuantas veces más.}

\subsection{Resumen de la sección}

Retrieval aporta a reasoning hechos externos, casos y principios abstractos; reasoning convierte esa evidencia en respuestas. Analogical prompting, step-back y deep research pueden verse como un espectro continuo que va de una analogía única a una investigación de varias rondas. La calidad del sistema debe descomponerse en tasa de recuperación, confiabilidad de la evidencia, corrección del razonamiento y fidelidad de las citas.
